%% Согласно ГОСТ Р 7.0.11-2011:
%% 5.3.3 В заключении диссертации излагают итоги выполненного исследования, рекомендации, перспективы дальнейшей разработки темы.
%% 9.2.3 В заключении автореферата диссертации излагают итоги данного исследования, рекомендации и перспективы дальнейшей разработки темы.
\begin{enumerate}
  \item Проанализирована задача автоматической классификации электронной почты
        с~учётом персонализированных пользовательских классов, приватности
        данных и~дрейфа тем. На~этой основе сформулирована многоклассовая
        постановка с~порогом уверенности и~служебной категорией
        \texttt{other} для отказа от~низкоуверенной классификации.
  \item Разработана локальная архитектура обработки почтового архива:
        синхронизация и~нормализация писем, построение многоязычных
        эмбеддингов, кластеризация исторической выборки, ручная проверка
        кандидатов в~классы, построение прототипов, применение классификатора
        и~контроль обновления модели.
  \item Реализован прототип системы, не~передающий текст писем во~внешние
        сервисы. В~эксперименте использованы эмбеддинги \texttt{BAAI/bge-m3},
        плотностная кластеризация HDBSCAN с~подбором параметров и~прототипный
        классификатор поверх векторных представлений.
  \item Экспериментальная проверка на~обезличенной выборке из~7107 писем
        подтвердила выполнение основных критериев успеха: при~рабочем пороге
        \(\tau=0{,}55\) классификатор достиг macro-F1~\(0{,}835\),
        weighted-F1~\(0{,}847\) и~покрытия~\(0{,}996\) по~восьми
        пользовательским классам.
\end{enumerate}
